% Beamer-spezifische Konfiguration
% Beamer-specific configuration

%%%%%%%%%%%%%%%%%%%%%%%%%%%%%%%%%%%%%%%%%%%%%%%%%%%

% Einstellungen für das Dokument

%%%%%%%%%%%%%%%%%%%%%%%%%%%%%%%%%%%%%%%%%%%%%%%%%%%

% Eine neue Farbe für keywords im Code und für links
\definecolor{codeblue}{HTML}{4178f2}

%Farbe für Textblocktitel
\definecolor{lieblingsfarbe}{cmyk}{79,0.11,1,0.25}

% Entfernt den roten Rahmen um klickbare Links und macht Links blau
\hypersetup{hidelinks, colorlinks=true, allcolors=codeblue}

% Stildefinition für Aufzählungspunkte
\setlist[itemize]{label=$\bullet$, leftmargin=*, nosep}

% Vordefinierte Umgebungen für Theoreme/Lemmas/Definitionen
\theoremstyle{definition}
\newtheorem{satz}[theorem]{Satz}
\newtheorem{fundamentalsatz}[theorem]{Fundamentalsatz}
\newtheorem{intuition}[theorem]{Intuition}
\newtheorem{intuitionUndBeispiel}[theorem]{Intuition und Beispiel}
\newtheorem{beispiel}[theorem]{Beispiel}
\newtheorem{korollar}[theorem]{Korollar}
\newtheorem{bemerkung}[theorem]{Bemerkung}
\newtheorem{vorbemerkung}[theorem]{Vorbemerkung}

% Definition der Trennfolien
\newcommand{\trennfolie}[1]{
	\begin{frame}[plain]
		\centering
		\vfill
		\Huge \textbf{#1}
		\vfill
	\end{frame}
}

% Verwende das enumitem-Paket, um das Aussehen der Aufzählungspunkte anzupassen
\setbeamertemplate{items}[default]
\setbeamertemplate{enumerate items}[circle]

% Beweisumgebung im gleichen Stil wie die anderen Umgebungen
\newenvironment{beweis}[1][\textbf{Beweis}]{%
	\begin{proof}[#1]
	}{%
	\end{proof}
}

% Stil für TikZ - Grafiken
\tikzset{
	heapnode/.style={draw, circle, inner sep=0pt, minimum size=6mm},
	heaparrow/.style={-latex, thick},
}

% Ersetze U+202F durch ein normales Leerzeichen
\newunicodechar{ }{ }

% Ersetze U+0308 durch das LaTeX-Äquivalent für das Diaeresis (z.B., \"{o} für ö)
\newunicodechar{̈}{\"}

%%%%%%%%%%%%%%%%%%%%%%%%%%%%%%%%%%%%%%%%%%%%%%%%%%%

% Definiere Sprache 'pseudocode' für Codelistings
\lstdefinelanguage{pseudocode}{
	keywords={
		struct,
		null,
		function,
		error,
		nil,
		if,
		then,
		else,
		for,
		while,
		do,
		in,
		True,
		False,
		Array,
		set,
		repeat,
		break,
		continue,
		Eingabe,
		Ausgabe,
		return,
		print,
		exit
	},
	sensitive=false,
	delim=[l][keywordstyle]{:},
	comment=[l][commentstyle]{\#}
}

% Generelle Einstellungen für Listings
\lstset{%
	frame=single,
	framesep=0mm,
	framexleftmargin=7mm,
	xleftmargin=8mm,
	framerule=1.5pt,
	rulecolor=\color{black},
	backgroundcolor=\color{white},
	basicstyle=\ttfamily,
	keywordstyle=\bfseries\color{codeblue},
	commentstyle=\color{gray},
	numbers=left,
	stepnumber=1,
	numbersep=1mm,
	numberstyle=\sffamily\color{gray!80!black}\footnotesize,
	numberblanklines=true,
	escapeinside={\%*}{*)},
	inputencoding=utf8
}

% Definiere einige Stile für TikZ - Knoten
\tikzset{
	long node/.style={draw, rounded corners, minimum width=7.5cm, minimum height=1cm},
	short node/.style={draw, rounded corners, minimum width=2.5cm, minimum height=1cm},
	square node/.style={draw, rounded corners, minimum width=2.5cm, minimum height=2.5cm},
	label node/.style={font=\small, align=center}
}

% Abkürzungen ("Makros") können mit \newcommand definiert werden
% Symbole für die gängigen Mengen
\newcommand{\N}{\mathbb{N}}
\newcommand{\Z}{\mathbb{Z}}
\newcommand{\Zp}{\Z_{\geq 0}}
\newcommand{\Q}{\mathbb{Q}}
\newcommand{\Qp}{\Q_{\geq 0}}
\newcommand{\R}{\mathbb{R}}
\newcommand{\Rp}{\R_{\geq 0}}
\newcommand{\Rmn}{\R^{m\times n}}

% Hübscheres Epsilon
\newcommand{\eps}{\varepsilon}

% Einführung von TikZ - Bibliotheken
\usetikzlibrary{positioning,fit,shapes.geometric}

%%%%%%%%%%%%%%%%%%%%%%%%%%%%%%%%%%%%%%%%%%%%%%%%%%%
% Metadaten: Autor, Titel, etc. (Beamer)
\author{Yavuzâlp Dal}
\title{Einf\"uhrung in die Kerndichtesch\"atzung}
\institute{Hochschule Düsseldorf}
\date{\today}
%%%%%%%%%%%%%%%%%%%%%%%%%%%%%%%%%%%%%%%%%%%%%%%%%%%
